% Bagian: counterfactuals
% Bab: introduction
% Subbagian: material-conditional

\documentclass[../../../include/open-logic-section]{subfiles}

\begin{document}

\olfileid[id]{cnt}{int}{mat}

\olsection{Kondisional Material}

Dalam bentuknya yang paling sederhana, kondisional dalam bahasa Indonesia
adalah kalimat berbentuk ``Jika \dots, maka \dots,'' dengan bagian-bagian
yang ditunjukkan oleh \dots{} itu sendiri merupakan kalimat, seperti ``Jika
kepala pelayan yang melakukannya, maka tukang kebun tidak bersalah.'' Dalam
kuliah logika pengantar, kita belajar menyimbolkan kondisional dengan
menggunakan konektif~$\lif$: bagian-bagian yang ditunjukkan oleh \dots
disimbolkan, misalnya, dengan !!{formula}~$!A$ dan~$!B$, lalu seluruh
kondisional disimbolkan dengan $!A \lif !B$.

Konektif~$\lif$ bersifat \emph{fungsional-kebenaran}, yakni nilai kebenaran---
$\True$ atau $\False$---dari $!A \lif !B$ ditentukan oleh nilai kebenaran
$!A$ dan~$!B$: $!A \lif !B$ benar jika dan hanya jika $!A$ salah atau $!B$
benar, dan salah dalam keadaan lainnya. Secara relatif terhadap suatu
penetapan nilai kebenaran~$\pAssign v$, kita mendefinisikan
$\pSat{v}{!A \lif !B}$ jika dan hanya jika $\pSat/{v}{!A}$ atau
$\pSat{v}{!B}$. Konektif~$\lif$ dengan semantik ini disebut
\emph{kondisional material}.

Definisi ini menghasilkan sejumlah fakta logis dasar. Pertama-tama, teorema
deduksi berlaku bagi kondisional material:
\begin{align}
  \text{Jika } \Gamma, !A \Entails !B \text{, maka }
  \Gamma \Entails !A \lif !B
\end{align}
Kondisional ini bersifat fungsional-kebenaran: $!A \lif !B$ dan
$\lnot !A \lor !B$ ekuivalen:
\begin{align}
  !A \lif !B & \Entails \lnot !A \lor !B\\
  \lnot !A \lor !B & \Entails !A \lif !B
  \intertext{Kondisional material diikutkan oleh konsekuennya dan oleh
    negasi antesedennya:}
  !B & \Entails !A \lif !B\\
  \lnot !A & \Entails !A \lif !B
  \intertext{Kondisional material yang salah ekuivalen dengan konjungsi
    antesedennya dan negasi konsekuennya: jika $!A \lif !B$ salah, maka
    $!A \land \lnot !B$ benar, dan sebaliknya:}
  \lnot(!A \lif !B) & \Entails !A \land \lnot !B\\
  !A \land \lnot !B & \Entails \lnot(!A \lif !B)
  \intertext{Kondisional material mendukung modus ponens:}
  !A, !A \lif !B & \Entails !B
  \intertext{Kondisional material memungkinkan aglomerasi:}
  !A \lif !B,  !A \lif !C & \Entails !A \lif (!B \land !C)
  \intertext{Kita selalu dapat memperkuat anteseden; yakni, kondisional
    tersebut \emph{monoton}:}
  !A \lif !B & \Entails (!A \land !C) \lif !B
  \intertext{Kondisional material bersifat transitif; yakni, aturan rantai
    valid:}
  !A \lif !B, !B \lif !C & \Entails !A \lif !C
  \intertext{Kondisional material ekuivalen dengan kontraposisinya:}
  !A \lif !B & \Entails \lnot !B \lif \lnot !A\\
  \lnot !B \lif \lnot !A & \Entails !A \lif !B
\end{align}

Semua inferensi ini berguna dan tidak bermasalah dalam penalaran matematis.
Namun, literatur filsafat dan linguistik penuh dengan contoh tandingan yang
dikemukakan terhadap inferensi-inferensi sepadan dalam konteks
nonmatematis. Contoh-contoh tersebut menyiratkan bahwa kondisional
material~$\lif$ bukan---atau setidaknya tidak selalu merupakan---konektif
yang tepat untuk digunakan ketika menyimbolkan pernyataan bahasa alamiah
berbentuk ``jika \dots, maka \dots.''

\end{document}
